% problem_id: S001-lemma-genus-nonnegative
% source_id: S001 (Stacks Project)
% source_file: models.tex
% statement_locator: models.tex lines 694--698
% proof_locator: models.tex lines 700--843
% env: lemma
% id_note: label 在本来源内唯一
% pairing: tight (verified)
% status: complete (statement + author proof verbatim + reason + locators)
%
\begin{lemma}
\label{lemma-genus-nonnegative}
If $n, m_i, a_{ij}, w_i, g_i$ is a minimal numerical type
with $n > 1$, then $g \geq g_{top}$.
\end{lemma}

\begin{proof}
The reader who is only interested in the case of numerical types
associated to proper regular models can skip this proof as we will
reprove this in the geometric situation later.
We can write
$$
g_{top} = 1 - n + \frac{1}{2}\sum\nolimits_{a_{ij} > 0} 1 =
1 + \sum\nolimits_i (-1 +
\frac{1}{2}\sum\nolimits_{j \not = i,\ a_{ij} > 0} 1) 
$$
On the other hand, we have
\begin{align*}
g & =
1 + \sum m_i(w_i(g_i - 1) - \frac{1}{2} a_{ii}) \\
& =
1 + \sum m_iw_ig_i - \sum m_iw_i +
\frac{1}{2} \sum\nolimits_{i \not = j} a_{ij}m_j \\
& =
1 + \sum\nolimits_i
m_iw_i(-1 + g_i + \frac{1}{2} \sum\nolimits_{j \not = i} \frac{a_{ij}}{w_i})
\end{align*}
The first equality is the definition, the second equality uses that
$\sum a_{ij}m_j = 0$, and the last equality uses that
uses $a_{ij} = a_{ji}$ and switching order of
summation. Comparing with the formula for $g_{top}$ we conclude
that the lemma holds if
$$
\Psi_i =
m_iw_i(-1 + g_i + \frac{1}{2} \sum\nolimits_{j \not = i} \frac{a_{ij}}{w_i})
- (-1 + \frac{1}{2}\sum\nolimits_{j \not = i,\ a_{ij} > 0} 1)
$$
is $\geq 0$ for each $i$. However, this may not be the case.
Let us analyze for which indices we can have $\Psi_i < 0$.
First, observe that
$$
(-1 + g_i + \frac{1}{2}\sum\nolimits_{j \not = i} \frac{a_{ij}}{w_i}) \geq
(-1 + \frac{1}{2}\sum\nolimits_{j \not = i,\ a_{ij} > 0} 1)
$$
because $a_{ij}/w_i$ is a nonnegative integer. Since $m_iw_i$ is
a positive integer we conclude that $\Psi_i \geq 0$ as soon as
either $m_iw_i = 1$ or the left hand side of the inequality is $\geq 0$
which happens if $g_i > 0$, or $a_{ij} > 0$ for at least two indices $j$, or
if there is a $j$ with $a_{ij} > w_i$. Thus
$$
P = \{i : \Psi_i < 0\}
$$
is the set of indices $i$ such that $m_iw_i > 1$, $g_i = 0$,
$a_{ij} > 0$ for a unique $j$, and $a_{ij} = w_i$ for this $j$.
Moreover
$$
i \in P \Rightarrow \Psi_i = \frac{1}{2}(-m_iw_i + 1)
$$
The strategy of proof is to show that given $i \in P$ we can borrow a bit
from $\Psi_j$ where $j$ is the neighbour of $i$, i.e., $a_{ij} > 0$.
However, this won't quite work because $j$ may be an index with $\Psi_j = 0$.

\medskip\noindent
Consider the set
$$
Z = \{j : g_j = 0\text{ and }
j\text{ has exactly two neighbours }i, k\text{ with }
a_{ij} = w_j = a_{jk}\}
$$
For $j \in Z$ we have $\Psi_j = 0$. We will consider sequences
$M = (i, j_1, \ldots, j_s)$ where $s \geq 0$,
$i \in P$, $j_1, \ldots, j_s \in Z$, and
$a_{ij_1} > 0, a_{j_1j_2} > 0, \ldots, a_{j_{s - 1}j_s} > 0$.
If our numerical type consists of two indices which are in $P$
or more generally if our numerical type consists of
two indices which are in $P$ and all other indices in $Z$, then
$g_{top} = 0$ and we win by
Lemma \ref{lemma-non-irreducible-minimal-type-genus-at-least-one}.
We may and do discard these cases.

\medskip\noindent
Let $M = (i, j_1, \ldots, j_s)$ be a maximal sequence and let
$k$ be the second neighbour of $j_s$. (If $s = 0$, then $k$
is the unique neighbour of $i$.) By maximality $k \not \in Z$
and by what we just said $k \not \in P$. Observe that
$w_i = a_{ij_1} = w_{j_1} = a_{j_1j_2} = \ldots = w_{j_s} =
a_{j_sk}$. Looking at the definition
of a numerical type we see that
\begin{align*}
m_ia_{ii} + m_{j_1}w_i & = 0,\\
m_iw_i + m_{j_1}a_{j_1j_1} + m_{j_2}w_i & = 0,\\
\ldots & \ldots \\
m_{j_{s - 1}}w_i + m_{j_s}a_{j_sj_s} + m_kw_i & = 0
\end{align*}
The first equality implies $m_{j_1} \geq 2m_i$ because the
numerical type is minimal. Then the second equality implies
$m_{j_2} \geq 3m_i$, and so on. In any case, we conclude that
$m_k \geq 2m_i$ (including when $s = 0$).

\medskip\noindent
Let $k$ be an index such that we have a $t > 0$ and pairwise distinct
maximal sequences $M_1, \ldots, M_t$ as above, with
$M_b = (i_b, j_{b, 1}, \ldots, j_{b, s_b})$
such that $k$ is a neighbour of $j_{b, s_b}$ for $b = 1, \ldots, t$.
We will show that $\Phi_j + \sum_{b = 1, \ldots, t} \Phi_{i_b} \geq 0$.
This will finish the proof of the lemma by what we said above.
Let $M$ be the union of the indices occurring in $M_b$, $b = 1, \ldots, t$.
We write
$$
\Psi_k =
-\sum\nolimits_{b = 1, \ldots, t} \Psi_{i_b} + \Psi_k'
$$
where
\begin{align*}
\Psi_k' & =
m_kw_k\left(-1 + g_k +
\frac{1}{2} \sum\nolimits_{b = 1, \ldots t}
(\frac{a_{kj_{b, s_b}}}{w_k} - \frac{m_{i_b}w_{i_b}}{m_kw_k}) +
\frac{1}{2} \sum\nolimits_{l \not = k,\ l \not \in M}
\frac{a_{kl}}{w_k}
\right) \\
&
-\left(
-1 + \frac{1}{2}\sum\nolimits_{l \not = k,\ l \not \in M,\ a_{kl} > 0} 1
\right)
\end{align*}
Assume $\Psi_k' < 0$ to get a contradiction.
If the set $\{l : l \not = k,\ l \not \in M,\ a_{kl} > 0\}$ is empty,
then $\{1, \ldots, n\} = M \cup \{k\}$ and $g_{top} = 0$
because $e = n - 1$ in this case and the result holds by
Lemma \ref{lemma-non-irreducible-minimal-type-genus-at-least-one}.
Thus we may assume there is at least one such $l$ which contributes
$(1/2)a_{kl}/w_k \geq 1/2$ to the sum inside the first brackets.
For each $b = 1, \ldots, t$ we have
$$
\frac{a_{kj_{b, s_b}}}{w_k} - \frac{m_{i_b}w_{i_b}}{m_kw_k} =
\frac{w_{i_b}}{w_k}(1 - \frac{m_{i_b}}{m_k})
$$
This expression is $\geq \frac{1}{2}$ because $m_k \geq 2m_{i_b}$
by the previous paragraph and is $\geq 1$ if $w_k < w_{i_b}$.
It follows that $\Psi_k' < 0$ implies $g_k = 0$.
If $t \geq 2$ or $t = 1$ and $w_k < w_{i_1}$, then $\Psi_k' \geq 0$
(here we use the existence of an $l$ as shown above) which
is a contradiction too.
Thus $t = 1$ and $w_k = w_{i_1}$. If there at least two nonzero terms
in the sum over $l$ or if there is one such $k$ and $a_{kl} > w_k$, then
$\Psi_k' \geq 0$ as well. The final possibility is that $t = 1$ and
there is one $l$ with $a_{kl} = w_k$. This is disallowed as this would
mean $k \in Z$ contradicting the maximality of $M_1$.
\end{proof}
